\subsubsection{The \texorpdfstring{$\Sig_3$}{Sigma3} presentations of the sextic: elementary
fails, dressed closes}
\label{ped:sec:milestones:sigma3}

The dressed vertex \eqref{p:eq:hs6:dressed-vertex} carries $V_4^{-1}$ \emph{inside} a three-$\Sig_2$
coupling --- an unusual object. (In this and the next subsection we write $\Sig_2\equiv\Sig$
for the bilocal auxiliary, to keep the tower index explicit.) It is natural to ask whether one can
avoid the composite altogether by introducing the \emph{next} field of the auxiliary
tower early: a three-string auxiliary
\begin{equation}
	\Sig_3\in\Sym^3\Hilb_0,\qquad \Ngh(\Sig_3)=6,
	\label{ped:eq:Sigma3def}
\end{equation}
(the tower assignment $\Ngh(\Sig_k)=2k$; $\Sig_3$ is first \emph{needed} only at
$V_7=(3{,}2{,}2)$, but nothing forbids introducing it at $O(\gs^4)$). The hope: replace the
composite $W_6(\Sig_2^3)$ by \emph{elementary} vertices only --- a bilinear kinetic term
$V_6(\Sig_3,\Sig_3)$ and a cubic source $V_6(\Sig_3,\Sig_2,\Psi)$ --- so that the
$V_4^{-1}$ inside $W_6$ is traded for a $\Sig_3$-propagator $V_6^{-1}$. Ghost saturation is
consistent: $V_6(\Sig_3,\Sig_3)$ saturates at $6{+}6=12=2\cdot6$, and
$V_6(\Sig_3,\Sig_2,\Psi)$ at $6{+}4{+}2=12$. Both presentations below are machine-checked
on generic (non-$\Loo$) tensors in a finite-dimensional model (the scripts are listed in \cref{ch:finite-models}).

\subparagraph{The elementary presentation fails.} Take
\begin{equation}
	S_{\rm elem}
	= S_3
	+ \tfrac{\gs^4}{6!}\big[-V_6(\Sig_3,\Sig_3)+2\,V_6(\Sig_3,\Sig_2,\Psi)\big],
	\label{ped:eq:Selem}
\end{equation}
with $S_3$ the quintic action, $\caction[\le 4]$ plus \eqref{p:eq:hs5:S5tilde} (the $-1{:}2$ kinetic-to-source ratio mirrors
\eqref{p:eq:hs4:S4tilde}). Since $\Sig_3$ enters only quadratically$+$linearly, it can be eliminated
exactly (Gaussian): with $M:=V_6|_{\Sym^3}$ and $u(\Sig_2):=V_6(\cdot,\Sig_2,\Psi)$, the
elimination produces the $\Sig_3$-channel exchange
$\mathcal D(\Sig_2,\Sig_2)=u^\top M^{-1}u$, \emph{quadratic} in $\Sig_2$. Three facts kill
the hope:
\begin{enumerate}[label=(\roman*),leftmargin=*]
	\setlength{\itemsep}{0.4ex}
	\item \emph{On-shell, $\mathcal D$ collapses to the bare contact.} By $\Sym^3$
	completeness and total symmetry of $V_6$, the source collapses on the Veronese diagonal:
	$u(P)=M\,p_3$ with $p_3:=\Psi^{\otimes3}$, hence
	$\mathcal D(P,P)=p_3^\top M\,p_3=V_6(\Psi^6)$ --- a telescoping $V_6V_6^{-1}V_6=V_6$,
	the exact analogue of the quartic collapse \eqref{p:eq:hs4:telescoping}. The $\Sig_3$ sector supplies precisely the
	\emph{target} sextic contact, nothing else.
	\item \emph{The feedback is a $\Sig_2$-channel object and survives.} The quintic
	overshoot $\tfrac{\gs^4}{600}vV_4^{-1}v$ is built from $V_4,V_5$ \emph{only} --- it is
	exactly $V_6$-independent, so no $V_6$-built exchange can cancel it. Structurally:
	$\mathcal D$ is quadratic in $\Sig_2$, while the composite it would have to reproduce,
	$\mathcal C=V_5V_4^{-1}V_5$, is cubic in $\Sig_2$ --- different field degrees, different
	invariants. The joint on-shell mismatch $\tfrac{\gs^4}{600}vV_4^{-1}v$ remains for
	generic tensors.
	\item \emph{No repair exists within the content.} A cubic-in-$\Sig_2$ $\Sig_3$-source
	would need $V_7(\Sig_3,\Sig_2,\Sig_2)$ --- septic, hence $O(\gs^5)$, forbidden at this
	order; and a $\Sig_3$-source at order $\gs^k$ feeds an exchange at $\gs^{2k-4}$, so only
	$k=4$ contributes at $O(\gs^4)$ at all.
\end{enumerate}
The conclusion: \emph{the composite $\mathcal C$ is
irreducible} --- no elementary $\Sig_2\oplus\Sig_3$ presentation closes the $O(\gs^4)$
matching; the $V_4^{-1}\to V_6^{-1}$ trade is not an identity. (A new auxiliarity condition
appears alongside the one for $\Sig_2$: $\Sig_3$ is purely auxiliary iff $V_6$ is non-degenerate
on $\Sym^3\Hilb_0$.)

\subparagraph{The dressed presentation closes.} The cure is not to reproduce the \emph{bare} $V_6(\Psi^6)$ with $\Sig_3$ but
to pair $\Sig_3$ with the \emph{dressed effective sextic} --- the vertex that already
contains the lower-order feedback. Define the linear map
\begin{equation}
	\Phi:\ \Sym^3\Hilb_0\to\big(\Sym^2\Hilb_0\big)^*,
	\qquad
	\Phi(X):=V_5(\,\cdot\,,X)\big|_{\Sym^2},
	\label{ped:eq:Phimap}
\end{equation}
i.e.\ contract $V_5$ on three of its five slots with $X\in\Sym^3$ and read the remaining
two slots as a covector on $\Sym^2$ (well defined by total symmetry of $V_5$). Its
key property is the \emph{Veronese collapse}
\begin{equation}
	\Phi(\Psi^{\otimes3})=v
	\qquad\big(v=V_5(\cdot,\Psi^{\otimes2},\Psi)\text{, the quintic source covector}\big):
	\label{ped:eq:Phicollapse}
\end{equation}
on the diagonal, splitting three $\Psi$-slots as $2{+}1$ is immaterial. Then set
\begin{equation}
	\boxed{\;
	\tilde M:=M-\tfrac65\,N,
	\qquad
	N:=\Phi^\top V_4^{-1}\Phi,
	\qquad
	\tilde u(\Sig_2):=u(\Sig_2)-\tfrac65\,\Phi^\top V_4^{-1}\,v(\Sig_2),
	\;}
	\label{ped:eq:tildeMdef}
\end{equation}
with $v(\Sig_2):=V_5(\cdot,\Sig_2,\Psi)$ the off-shell extension of $v$. Here $M$,
$\tilde M$ are bilinear forms on $\Sym^3\Hilb_0$, while $u(\Sig_2)$, $\tilde u(\Sig_2)$
are \emph{covectors} on $\Sym^3\Hilb_0$ (linear functionals of the would-be $\Sig_3$
argument), themselves linear in $\Sig_2$; we write $\Mean{\tilde u(\Sig_2),\Sig_3}$ for
the evaluation. Explicitly, in an orthonormal basis $\{\Lambda^2_I\}$ of $\Sym^2$:
\begin{equation}
	\begin{aligned}
	\Mean{\tilde u(\Sig_2),\Sig_3}
	&= V_6(\Sig_3,\Sig_2,\Psi)
	- \tfrac65\sum_{I,J} V_5(\Lambda^2_I,\Sig_3)\,(V_4^{-1})_{IJ}\,V_5(\Lambda^2_J,\Sig_2,\Psi)
	\\&=: \tilde V_6(\Sig_3,\Sig_2,\Psi),
	\end{aligned}
	\label{ped:eq:utildeexplicit}
\end{equation}
the first term the bare source, the second the $\Phi$-map bridge: one $V_5$ eats the
three strands of $\Sig_3$, the other eats $\Sig_2$ and $\Psi$, and a $\Sig_2$-propagator
$V_4^{-1}$ joins them. The action is
\begin{equation}
	S_{\rm dr}
	= S_3
	+ \tfrac{\gs^4}{6!}\big[-\tilde M(\Sig_3,\Sig_3)+2\,\Mean{\tilde u(\Sig_2),\Sig_3}\big].
	\label{ped:eq:Sdr}
\end{equation}
Both $\tilde M$ and $N$ are $\Psi$-\emph{independent} (built from $V_4,V_5,V_6$ and fixed
bases alone) and symmetric; this will matter for the gauge law. Two identities, both direct
consequences of \eqref{ped:eq:Phicollapse}, do all the work:
\begin{equation}
	\tilde M(p_3,p_3)
	= V_6(\Psi^6)-\tfrac65\,vV_4^{-1}v
	= W_6(\Psi^6),
	\qquad
	\tilde u(P)=\tilde M(p_3,\cdot).
	\label{ped:eq:tildecollapse}
\end{equation}
The second is an equality of covectors on $\Sym^3$: $\tilde u(P)$ means \enquote{the
source covector evaluated on the $\Sig_2$-saddle $\Sig_2=P=\Psi^{\otimes2}$}, i.e.\
$\Mean{\tilde u(P),X}=\tilde V_6(X,\Psi^{\otimes2},\Psi)$ for every $X\in\Sym^3$, and the
identity says this functional coincides with pairing $X$ against $\Psi^{\otimes3}$
through the dressed kinetic form, $\tilde M(p_3,X)$. It is proved in one line from
\eqref{ped:eq:Phicollapse}: $u(P)=M(p_3,\cdot)$ is the bare collapse of the elementary presentation (total symmetry of
$V_6$: splitting six $\Psi$-slots as $3{+}2{+}1$ or $3{+}3$ is immaterial), and
$\Phi^\top V_4^{-1}v\big|_{\Sig_2=P}=N(p_3,\cdot)$ by $\Phi(p_3)=v$. Consequently the
dressed saddle is \emph{exactly} $\Sig_3^*=\Psi^{\otimes3}$ on-shell, and the resulting
$\Sig_3$-contact is $\tfrac{\gs^4}{6!}\tilde M(p_3,p_3)=\tfrac{\gs^4}{6!}W_6(\Psi^6)$ ---
the geometric sextic \emph{and} the feedback subtraction, in one Gaussian step. The
matching closes exactly as in \cref{p:sec:hs:sextic}, with the composite
$-\tfrac65\mathcal C$ now living \emph{inside the $\Sig_3$-kinetic pairing} instead of
inside a three-$\Sig_2$ vertex: the map $\Phi$ is the relocation mechanism, converting the
$\Sig_2$-channel feedback $vV_4^{-1}v$ (a $\Sym^2$-channel object) into a $\Sym^3$-kinetic
correction $N$. The composite is not eliminated --- the elementary presentation shows it cannot be ---
it is \emph{repackaged} where it does the least harm: in a quadratic kernel, the natural
home of a propagator-like object, rather than in a cubic vertex.

\subparagraph{The complete cubic action through $O(\gs^4)$, both presentations.} For
reference:
\begin{equation}
	\boxed{
	\begin{aligned}
	S^{\rm min}_4
	&= \tfrac12V_2(\Psi^2)+\tfrac{\gs}{3!}V_3(\Psi^3)
	-\tfrac{\gs^2}{4!}V_4(\Sig_2,\Sig_2)+\tfrac{2\gs^2}{4!}V_4(\Sig_2,\Psi^{\otimes2})
	\\&\quad
	+\tfrac{\gs^3}{5!}V_5(\Sig_2,\Sig_2,\Psi)
	+\tfrac{\gs^4}{6!}\,W_6(\Sig_2,\Sig_2,\Sig_2),
	\\[1ex]
	S^{\rm dr}_4
	&= \tfrac12V_2(\Psi^2)+\tfrac{\gs}{3!}V_3(\Psi^3)
	-\tfrac{\gs^2}{4!}V_4(\Sig_2,\Sig_2)+\tfrac{2\gs^2}{4!}V_4(\Sig_2,\Psi^{\otimes2})
	\\&\quad
	+\tfrac{\gs^3}{5!}V_5(\Sig_2,\Sig_2,\Psi)
	-\tfrac{\gs^4}{6!}\tilde M(\Sig_3,\Sig_3)+\tfrac{2\gs^4}{6!}\Mean{\tilde u(\Sig_2),\Sig_3}.
	\end{aligned}
	}
	\label{ped:eq:S4both}
\end{equation}
Both are cubic in the tower fields $(\Psi,\Sig_2,\Sig_3)$, agree exactly with
$S^{(6)}_{\rm eff}$ on-shell up to $O(\gs^5)$ feedback, and carry the homogeneous gauge
laws of \cref{sec:order:sextic:both-gauge}. The minimal presentation is more economical (one field
fewer); the dressed presentation is structurally uniform --- every order adds one field and
one HS pair (kinetic $+$ source), which is what generalises
(\cref{ch:all-orders}).

\subparagraph{Equations of motion and the auxiliarity chain.} In both presentations the
$\Sig_3$-EOM is Gaussian --- $\Sig_3$ enters bilinearly $+$ linearly --- with exact saddle
$\Sig_3^*=\tilde M^{-1}\tilde u(\Sig_2)$ (dressed; bare with $M,u$), collapsing to
$\Psi^{\otimes3}$ on-shell by \eqref{ped:eq:tildecollapse}. The $\Sig_2$-EOM is the
\emph{implicit} saddle equation ($\tilde X_6=0$ with the $\Sig_3$-exchange
folded into the Hessian): the $\Sig_2$-saddle is a perturbative fixed point
$\Sig_2^*(\Psi)=\Psi^{\otimes2}+O(\gs)$, Neumann-solvable at small $\gs$. The auxiliarity
conditions chain up: $\Sig_2$ is purely auxiliary iff $V_4|_{\Sym^2}$ is non-degenerate
(\cref{p:sec:hs:quartic}); $\Sig_3$ iff the relevant kinetic kernel is --- $M=V_6|_{\Sym^3}$
(elementary) or $\tilde M$ (dressed; this needs \emph{both} $V_4$ and $V_6$ well behaved,
tying the two conditions together). Reality extends verbatim from
\cref{p:eq:hs4:reality} (in the sign convention of the translation box): $\Sig_3^\ddagger=\Sig_3$, inherited from
$(\Psi^{\otimes3})^\ddagger=\Psi^{\otimes3}$, makes every new term real.

